\subsection{GRADUATIONS}

Let \(\Delta\) be a commutative monoid, written additively. Let \(\phi_0\) denote a mapping of X into \(\Delta\) and \(\phi\) the homomorphism of the free magma M(X) into \(\Delta\) which extends \(\phi_0\) . For all \(\delta \in \Delta\) , let \(\operatorname{Lib}^{\delta}(X)\) be the submodule of \(\operatorname{Lib}(X)\) with basis the subset \(\phi^{-1}(\delta)\) of M(X). The family \((\operatorname{Lib}^{\delta}(X))_{\delta \in \Delta}\) is a graduation of the algebra \(\operatorname{Lib}(X)\) , that is

\begin{enumerate}
\def\labelenumi{(\arabic{enumi})}
\setcounter{enumi}{7}
\tightlist
\item
\end{enumerate}

\[
\operatorname{Lib} (\mathrm{X}) = \bigoplus_ {\delta \in \Delta} \operatorname{Lib} ^ {\delta} (\mathrm{X})\tag{8}
\]

\[
\operatorname{Lib} ^ {\delta} (\mathrm{X}). \operatorname{Lib} ^ {\delta^ {\prime}} (\mathrm{X}) \subset \operatorname{Lib} ^ {\delta + \delta^ {\prime}} (\mathrm{X}) \quad \text { for } \delta , \delta^ {\prime} \text { in } \Delta\tag{9}
\]

(Algebra, Chapter III, § 3, no. 1, Example 3).

Lemma 1. The ideal \(\mathfrak{a}\) of Definition 1 is graded.

For \(a, b\) in \(\operatorname{Lib}(\mathbf{X})\) , let \(\mathbf{B}(a, b) = a.b + b.a\) . The formulae

\begin{enumerate}
\def\labelenumi{(\arabic{enumi})}
\setcounter{enumi}{9}
\tightlist
\item
\end{enumerate}

\[
\mathrm{B} (a, b) = \mathrm{Q} (a + b) - \mathrm{Q} (a) - \mathrm{Q} (b)\tag{10}
\]

\[
Q \left(\lambda_ {1}. w _ {1} + \dots + \lambda_ {n}. w _ {n}\right) = \sum_ {i} \lambda_ {i} ^ {2} Q \left(w _ {i}\right) + \sum_ {i <   j} \lambda_ {i} \lambda_ {j} B \left(w _ {i}, w _ {j}\right)\tag{11}
\]

for \(w_1, \ldots, w_n\) in M(X) and \(\lambda_1, \ldots, \lambda_n\) in K, show that the families \((\mathbf{Q}(a))_{a \in \mathrm{Lib}(\mathbf{X})}\) and \((\mathbf{Q}(w), \mathbf{B}(w, w'))_{w, w' \in \mathrm{M}(\mathbf{X})}\) generate the same submodule of Lib(X). As J is trilinear the ideal a is generated by the homogeneous elements \(\mathbf{Q}(w), \mathbf{B}(w, w')\) and \(\mathbf{J}(w, w', w'')\) , where \(w, w', w''\) are in M(X), and hence is graded (Algebra, Chapter III, § 3, no. 3, Proposition 1).

Let the Lie algebra \(\mathrm{L}(\mathbf{X}) = \mathrm{Lib}(\mathbf{X}) / a\) be given the quotient graduation. The homogeneous component of \(\mathrm{L}(\mathbf{X})\) of degree \(\delta\) is denoted by \(\mathrm{L}^{\delta}(\mathbf{X})\) ; it is the submodule of \(\mathrm{L}(\mathbf{X})\) generated by the images of the elements \(w \in \mathrm{M}(\mathbf{X})\) such that \(\phi(w) = \delta\) .

We shall make special use of the following two cases:

\begin{enumerate}
\def\labelenumi{(\alph{enumi})}
\tightlist
\item
  Total graduation: we take \(\Delta = \mathbf{N}\) and \(\phi_0(x) = 1\) for all \(x\in \mathbf{X}\) , whence \(\phi (w) = l(w)\) for \(w\) in \(\mathrm{M}(\mathrm{X})\) . The K-module \(\mathrm{L}^n (\mathrm{X})\) is generated by the images of the elements of length \(n\) in \(\mathrm{M}(\mathrm{X})\) , which we shall call alternants of degree \(n\) . We shall see later that the module \(\mathrm{L}^n (\mathrm{X})\) is free and admits a basis consisting of alternants of degree \(n\) (no. 11, Theorem 1). Then \(\mathrm{L}(\mathrm{X}) = \bigoplus_{n\geq 1}\mathrm{L}^n (\mathrm{X})\) and \(\mathrm{L}^1 (\mathrm{X})\) admits X as basis (no. 2, Corollary 1 to Proposition 1). By the construction of \(\mathrm{M}(\mathrm{X})\) ,
\end{enumerate}

\[
\mathrm{L} ^ {n} (\mathrm{X}) = \sum_ {p = 1} ^ {n - 1} [ \mathrm{L} ^ {p} (\mathrm{X}), \mathrm{L} ^ {n - p} (\mathrm{X}) ]\tag{12}
\]

and in particular

\[
[ \mathrm{L} ^ {m} (\mathrm{X}), \mathrm{L} ^ {n} (\mathrm{X}) ] \subset \mathrm{L} ^ {m + n} (\mathrm{X}).\tag{13}
\]

\begin{enumerate}
\def\labelenumi{(\alph{enumi})}
\setcounter{enumi}{1}
\tightlist
\item
  Multigraduation: we take \(\Delta\) to be the free commutative monoid \(\mathbf{N}^{(\mathbf{X})}\) constructed on X. The mapping \(\phi_0\) of X into \(\Delta\) is defined by \((\phi_0(x))(x') = \delta_{xx'}\) , where \(\delta_{xx'}\) is the Kronecker symbol. For \(w \in M(X)\) and \(x \in X\) , the integer \((\phi(w))(x)\) is ``the number of occurrences of the letter \(x\) in \(w\)''. For \(\alpha\) in \(\mathbf{N}^{(\mathbf{X})}\) , we write \(|\alpha| = \sum_{x \in X} \alpha(x)\) , whence \(|\phi(w)| = l(w)\) for all \(w\) in M(X). It follows that
\end{enumerate}

\[
\mathbf {L} ^ {n} (\mathbf {X}) = \bigoplus_ {| \alpha | = n} \mathbf {L} ^ {\alpha} (\mathbf {X});\tag{14}
\]

obviously

\[
[ \mathrm{L} ^ {\alpha} (\mathrm{X}), \mathrm{L} ^ {\beta} (\mathrm{X}) ] \subset \mathrm{L} ^ {\alpha + \beta} (\mathrm{X}) \quad \text { for } \alpha , \beta \text { in } \mathbf {N} ^ {(\mathrm{X})}.\tag{15}
\]

PROPOSITION 4. Let \(S\) be a subset of \(X\) . If \(\mathbf{N}^{(\mathbb{S})}\) is identified with its canonical image in \(\mathbf{N}^{(\mathbf{X})}\) (Algebra, Chapter I, § 7, no. 7), then \(L(S) = \sum_{\alpha \in \mathbf{N}^{(\mathbb{S})}} L^{\alpha}(X)\) . Further, for all \(\alpha \in \mathbf{N}^{(\mathbb{S})}\) , the homogeneous component of degree \(\alpha\) under the multigraduation on \(L(S)\) is equal to \(L^{\alpha}(X)\) .

 Let \(\alpha \in \mathbf{N}^{(S)}\) . The module \(\mathrm{L}^{\alpha}(\mathrm{S})\) is generated by the images in \(\mathrm{L}(\mathrm{X})\) of the elements \(w\) in \(\mathrm{M}(\mathrm{S})\) such that \(\phi(w) = \alpha\) , that is (Algebra, § 7, no. 9, formulae (23) and (24)) the set of \(w\) in \(\mathrm{M}(\mathrm{X})\) such that \(\phi(w) = \alpha\) . Hence \(\mathrm{L}^{\alpha}(\mathrm{S}) = \mathrm{L}^{\alpha}(\mathrm{X})\) . The proposition follows from this and the relation \(\mathrm{L}(\mathrm{S}) = \sum_{\alpha \in \mathbf{N}^{(S)}}\mathrm{L}^{\alpha}(\mathrm{S})\) .

COROLLARY. For every family \((\mathbf{S}_i)_{i\in I}\) of subsets of \(\mathbf{X}\) ,

\[
\mathrm{L} \left(\bigcap_ {i \in I} S _ {i}\right) = \bigcap_ {i \in I} \mathrm{L} (S _ {i}).\tag{16}
\]

This follows from Proposition 4 and the obvious formula

\[
\mathbf {N} ^ {(S)} = \bigcap_ {i \in I} \mathbf {N} ^ {(S _ {i})}\tag{17}
\]

where we have written \(S = \bigcap_{i \in I} S_i\) .

